\section{Representation variety}
We refer to \cite{maret_note_nodate, casimiroCharacterVarietiesTheir2026} for further details.


Let $\Gamma$ be a finitely generated group and $G$ be a (complex or real) Lie Group.

\begin{definition}
	Given a finitely generated group $\Gamma$ and a Lie Group, we define the representation variety denoted by $\operatorname{Hom}(\Gamma,G)$ as 	
	$$ \operatorname{Hom}(\Gamma,G) = \{\rho: \Gamma \to G | \rho \text{ is a group homomorphism}\} $$	
\end{definition}
the elements $\varphi \in \Hom(\Gamma,G)$ are called representations.

The representation variety carries the subspace topology induced from the compact-open topology on the space of continuous maps $\Gamma \to G$ where $\Gamma$ is endowed with the discrete topology. A subbasis for this topology is given by the sets
$$ V(K,U) = \{f: \Gamma \to G | K \subset \Gamma \text{ finite}, U \subset \text{ open}, f(K)\subset U \}$$


It is usual to think of the representation variety as the space 
$$ X(\Gamma, G) = \{\rho(\gamma_{1}), \dots, \rho(\gamma_n) | \rho \in \operatorname{Hom}(\Gamma, G)   \} \subset G^n$$

Indeed, we will see that this space is homeomorphic to the representation variety

\begin{proposition}
 If $G$ is a Lie group endowed with an analityc atlas then $X(\Gamma,G)$ is an analytic subvariety of $G^n$, homeomorphic to $\Hom(\Gamma,G)$. In particular $\Hom(\Gamma,G)$ inherits a natural structure of analityc variety wich is independent of the choice of generators of $\Gamma$.  	
\end{proposition}

\begin{proof}
	Let $\{\gamma_{1}, \dots, \gamma_n\}$ be a set of generators of $\gamma$, and let $r_i$ denote the set of relations for these generators. Recall that every group homomorphis is defined by the image of its generators, and those images have to preserve the relations between the generators. Each relation can be thought as a map $r_i:\Gamma^n \to \Gamma$, since this map is just a composition of inversions and multiplications, then each $r_i$ is analytic. Then $X(\Gamma,G)$ is precisely the subset of $(g_{1},\dots,g_n) \in G^n$ defined by the conditions $r_i(g_{1},\dots,g_n)=1$. It is clear that we have an isomorphism
	\[
		\Psi: \Hom(\Gamma,G) \to X(\Gamma,G), \quad \rho \to (\rho(\gamma_{1}),\dots,\rho(\gamma_n))
	\]
To see that $\Psi$ is a homeomorphism, consider $U_{1},\dots,U_n$ open sets in $G$. Then 
\[
	\Psi^{-1}(X(\Gamma,G)\cap \left(U_1 \times \dots \times U_n\right)) = \Hom(\Gamma,G) \cap \bigcap_{i=1}^n V \left(\{\gamma_i)\}, U_i \right)
\]
This proves $\Psi$ is continuous. Now for its inverse, take $K \subset \Gamma$ finite and $U \subset G$ open. Any $\gamma \in K$ can be expressed as a word in the generators, and hence it determines an analytic map $\gamma:G^n \to G $, hence
\[
	\Psi(\Hom(\Gamma,G)\cap V(K,U)) = X(\Gamma,G) \cap \bigcap_{\gamma \in K} \gamma ^{-1}(U)
\]
Therefore, $\Psi$ is a homeomorphism.

If $(\gamma'_{1},\dots,\gamma'_n)$ is another set of generators of $\Gamma$ and $X'(\Gamma,G)$ is the associated space, then the map from $X(\Gamma,G)$ to $X(\Gamma,G)$ defined as the composition 
\[
	X(\Gamma,G) \to \Hom(\Gamma,G) \to X'(\Gamma,G)
\]is an isomorphism of analytic varieties. Indeed, the map sends $(\rho(\gamma_{1}),\dots,\rho(\gamma_n))$ to $(\rho(\gamma_{1}'),\dots,\rho(\gamma_n))$ . Now, since $\gamma_i'$ is a word in the generators $\gamma_{1},\dots,\gamma_n$, it follows that $\rho(\gamma_i')$ is a word in $\rho(\gamma_{1}),\dots,\rho(\gamma_n)$. This shows that the map is analytic because word maps are analytic by assumption on $G$.

\end{proof}
\begin{proposition}
	Assume that $G$ has the structure of a real or complex algebraic group, then $X(\Gamma,G)$ is an algebraic subset of $G^n$. In particular,$\Hom(\Gamma,G) $has a natural structure of real or complex algebraic variety and the structure does not depend on the choice of generators of $\Gamma$.
\end{proposition}

\begin{proof}
	The proof is analogous to the previous one.
\end{proof}

The representation variety $\Hom(\Gamma,G)$, has two natural symmetries given by the right action of the group $\operatorname{Aut}(\Gamma)$ of automorphisms of $\Gamma$ by pre-composition and the left action of $\operatorname{Aut}(G)$ by post-composition. Both actions preserve the analytic or algebraic structure of $\Hom(\Gamma,G)$ as a consequence of the next proposition
\begin{proposition}
	Let $\Gamma, \Gamma_{1},\Gamma_{2}$ be finitely generated groups, and let $G,G_{1},G_{2}$ be Lie or algebraic groups.
	\begin{enumerate}
	   \item If $f:\Gamma_{1} \to \Gamma_{2}$ is a homomorphism, then the induced map $ f^* : \Hom(\Gamma_{2}, G) \to \Hom(\Gamma_{1}, G),\quad \rho \mapsto \rho \circ f$  is analytic (or regular). 
	   \item If $\phi:G_{1}\to G_{2}$ is a homomorphism, then the induced map $ \phi^* : \Hom(\Gamma, G_{1}) \to \Hom(\Gamma, G_{2}), \rho \mapsto \phi \circ \rho$ is analytic (or regular).		   
	\end{enumerate}
\end{proposition}
\begin{proof}
	See \cite[Lemma 1.2.7]{maret_note_nodate}
\end{proof} 

Among $\operatorname{Aut(G)}$ normal subgroups, one is particularly important, the subgroup of inner automorphisms $\operatorname{Inn}(G)$ defined by conjugation with elements of $G$. In  particular, $\Inn(G)=G/Z(G)$. The action of $\operatorname{Inn}(G)$ on $\Hom(\Gamma, G)$ is fundamental in many geometric situations: for example, holonomy representations of surface structures are only well-defined up to conjugation by G. This motivates the study of the quotient $ \operatorname{Hom}(\Gamma, G)/\operatorname{Inn}(G)$,




\subsection{Characterizations of the orbits of the $\operatorname{Inn}(G)$-action }

\begin{lemma}
	The $\operatorname{Inn}(G)$-action is free on the $\operatorname{Inn}(G)$-invariant subset that consists of all the  representations $\phi$ such that $Z(G) = Z(\phi)$. Where $Z(\phi)$ is the centralizer of the image of $\phi$.
\end{lemma}

\begin{proof}
	Let $\phi$ be a representation satisfying the condition above. Let $\Psi_g \in \Inn(G)$, suppose that $\Psi_g(\phi) = \phi$, that is $g \phi g^{-1} = \phi$. Then $g \in Z(\phi) = Z(G)$, then $ \Psi_g(h)= g h g^{-1}=h$ for every $h \in G$, therefore $\Psi_g =e$
\end{proof}


\begin{proposition}[\cite{goldmanSymplecticNatureFundamental1984}]
	The action of $\Inn(G)$ on $\Hom(\Gamma,G)$ is locally free at $\phi$ if and only  if  $\operatorname{dim} Z(G)= \dim Z(\phi)$.
\end{proposition}
\begin{proof}
	\cite[Proposition 8]{casimiroCharacterVarietiesTheir2026}.
\end{proof}

The latter proposition motivates the following definition.

\begin{definition}
	A representation $\phi \in \Hom(\Gamma,G)$  is called regular if  $\dim Z(G) = \dim Z(\phi)$ and very regular if $Z(G)= Z(\phi)$.
\end{definition}



We denote by $\Homreg(\Gamma,G)$, the $\Inn(G)$-invariant subspace of regular representations.  The $\Inn(G)$-invariant subspace of very regular representations is denoted by $\Homvreg(\Gamma,G)$.  



\begin{definition}
	A proper group action is a continuous group action of a topological group $G$ on topological space $X$ for wich the map \[
		G \times X \to X\times X, (g,x) \to (gx,x)
	\]
	is a proper map, i.e., inverse images of compact sets are compact.
\end{definition}


\begin{definition}[polystable representations]
\end{definition}


\begin{theorem}[Nagata]
\end{theorem}


\section{GIT quotient}



\section{Character variety}

\subsection{Betti character variety}
